\section{Scope, Tiers and Evidence Method}\label{sec:method}
\answers{how the evidence behind RQ2--RQ7 was collected}

\paragraph*{Tiers by method} We split harnesses by what can be observed, not by market share.
\emph{Tier~A} (open source, white-box): we read the source at a pinned commit and run the client in containers (e.g.\ Codex CLI, OpenCode, Gemini CLI, Qwen Code, Aider, goose, Cline, Hermes Agent).
\emph{Tier~B} (closed but runnable, black-box): we install the official client in a disposable container, record what it writes and which hosts it contacts, and drive it headless (Claude Code, GitHub Copilot CLI, Antigravity CLI before login). Grok Build is an A/B hybrid: its source is public under Apache-2.0\evD~\cite{grokbuild_repo}, but the public tree lags the npm release (1.0.45 vs 1.0.52) and injects endpoints at build time, so we read the source white-box and treat the shipped npm client as Tier~B in E3 and E7.
\emph{Tier~C} (closed, not reachable): documentation, advisories and third-party research only (e.g.\ Cursor, Kiro, Devin, Jules).
\emph{Tier~D} (runtimes and SDKs) is scored on a separate ladder and never ranked against products (the rubric files and figures group it with hosted services as ``tier~C''). The matrix covers \nHarnessesAB{} tier A/B products and \nHarnessesC{} runtimes and hosted services.

\paragraph*{Black-box method} For closed harnesses we install the official release, read its documentation and advisories, parse its configuration, run it under \texttt{harnesswatch observe} (Section~\ref{sec:tool}), probe configuration trust with synthetic repositories carrying hooks, MCP servers and settings, and inspect package metadata. We never decompile, deobfuscate or use leaked source, never decrypt TLS or store request bodies, and stay within vendor terms (otherwise the cell is \emph{N/A (terms)}). No vendor login was placed in any container: closed clients were driven only through documented custom-endpoint features pointing at a local model (Claude Code's \texttt{ANTHROPIC\_BASE\_URL}, although Anthropic states that it ``doesn't support routing Claude Code to non-Claude models through any gateway''\evD~\cite{cc_llm_gateway}, so our Claude Code runs test harness behaviour in an unsupported configuration, not vendor-supported performance; Copilot CLI's bring-your-own-key mode, which needs no GitHub authentication\evD~\cite{copilot_byok}, and Grok Build's custom models\evD~\cite{grokbuild_overview}). Antigravity has no documented local-model path and its terms restrict use with other products\evD~\cite{ag_terms}, so it was observed only before login.

\paragraph*{Evidence labels} Every factual claim about a harness carries one of three labels, shown as superscripts:
\evT~\emph{verified-by-test}, behaviour or state we observed by running the pinned client;
\evD~\emph{from-official-docs}, vendor documentation, vendor advisories, registry metadata or the official source code (reading source counts as documentation, not as a test);
\evTP~\emph{third-party-report}, a write-up by an independent researcher or journalist.
Unverified claims are excluded. Documentation and source facts live in the evidence ledger and \texttt{matrix/evidence/}; every \evT{} fact traces to a probe or bench file under \texttt{experiments/results/} or \texttt{bench/results/}.

\paragraph*{Who did the work} Evidence collection, both matrix ratings and incident collection and classification were performed by LLM research agents run and supervised by the author. The two matrix raters scored from the same extracted fact lists, and incidents had a single classifier; Section~\ref{sec:limits} discusses what this means for reliability.

\paragraph*{Kinds of finding} We separate a \emph{vulnerability} (a defect that violates the product's own security intent, usually with a CVE or GHSA), a \emph{design weakness} (a documented default or capability that enables harm without being a bug, e.g.\ headless mode running repository hooks) and a \emph{policy choice} (e.g.\ telemetry or retention defaults). Only the first is a bug.

\paragraph*{Ethics} Experiments ran in disposable containers on the author's machine with inert canaries and synthetic secrets; the public-repository study stores aggregates only. Potentially unknown issues went to private coordinated disclosure and are not described here. For packages we inspected registry metadata only; one fact obtained by listing strings in a binary was withdrawn as outside this rule.
